% !TEX root = ../main.tex
% Insert this file as an appendix section.
% Status labels requested for this proof audit: "Proved" means "已证".
% Every asserted step below is proved; no unproved step is used.

\section{A large-center-stepsize lower bound}
\label{app:q3-large-alpha}

This appendix proves only the large primal-center stepsize branch for
\(s=c_s/\epsilon\). We use the initialization, update order, output,
problem class, and first hitting time of
Definitions~\ref{def:splalm}, \ref{def:problem-class}, and
\ref{def:complexity-measure}. The instance has the radial form of
Example~\ref{ex:rotation}, but its frequency and amplitude are chosen
differently. Neither the hypothesis \(\epsilon(\rho+s)\le c_1\) nor
\(\beta>0\) is imposed. In particular, the strict spectral-radius
conclusion of Lemma~\ref{lem:rot-spectral} is not invoked.

We identify \(\RR^2\) with \(\mathbb C\), as in the preliminaries.
Thus each complex coordinate below represents a real two-dimensional
block, its modulus is its Euclidean norm, and
\(\langle p,q\rangle=\operatorname{Re}(\overline p q)\).
All matrices acting on the four complex coordinates have real entries.

\subsection{Fixed constants and the statement}
\label{subsec:q3-large-constants}

\paragraph{Status: proved.}
Fix
\[
\rho>0,\qquad 0<\eta\le\gamma,\qquad
\tau>0,\qquad c_s>0,\qquad \Delta>0,
\qquad \chi:=\frac{\eta}{\gamma}\in(0,1].
\]
Choose the following explicit positive constants:
\begin{equation}
\begin{aligned}
\omega&:=\min\left\{
  \frac{\chi}{512},\,
  \frac{\sqrt{\eta\tau}}{2},\,
  \sqrt{\frac{\eta}{4(1+\rho/\tau)}}\right\},
&\delta&:=\frac{\omega^2}{64},\\
C_\lambda&:=\frac{400}{\chi^2},
&A_0&:=\min\left\{
  \frac14,\,\frac{\Delta}{4},\,
  \frac{1}{8(1+C_\lambda)+5\rho}\right\},\\
c_{\rm r}&:=\frac34 A_0\omega
  \min\left\{\frac{1}{\sqrt{\eta\tau}},\frac{1}{2\tau}\right\},
&B&:=\max\left\{1,\frac{\tau}{c_s\delta},\frac{2}{c_{\rm r}}\right\},\\
\epsilon_0&:=\min\left\{1,\frac{c_s}{2\tau},\frac{1}{2B}\right\}.
\end{aligned}
\label{eq:q3-large-fixed-constants}
\end{equation}
These constants depend only on
\((\rho,\gamma,\eta,\tau,c_s,\Delta)\). They do not depend on
\(\epsilon,\alpha,\beta\), or the instance dimensions.
The constant \(\delta\) here bounds \(d/\alpha\); it is not the
symbol \(\delta=1/s\) used in Appendix~\ref{app:upper-beta0}.

\begin{theorem}[Large primal-center stepsizes at \(s=c_s/\epsilon\)]
\label{thm:q3-large-alpha}
\textbf{Status: proved.}
Let the constants be given by~\eqref{eq:q3-large-fixed-constants}.
For every
\[
0<\epsilon<\epsilon_0,\qquad
\alpha\in(B\epsilon,1],\qquad \beta\in[0,1],\qquad
s:=\frac{c_s}{\epsilon},
\]
there exists an instance
\(I\in\mathcal I(\mathbf C_0,\Delta)\), with
\(\mathbf C_0=(1,1,0,1,1)\) and \(n=m=2\), such that SPLALM
started from zero satisfies
\[
R_k\ge c_{\rm r}\alpha>\epsilon
\qquad\text{for every }k\ge0.
\]
Consequently,
\[
T^{\mathrm{KKT}}_\epsilon(\theta(\epsilon),\alpha,\beta;I)
=+\infty,
\qquad
\theta(\epsilon)=(\rho,c_s/\epsilon,\gamma,\eta,\tau).
\]
\end{theorem}

The proof is completed in Step~5 below. Throughout the construction let
\[
d:=\frac{\tau}{s}=\frac{\tau\epsilon}{c_s}.
\]
The choice of \(B\) and the strict inequality \(\alpha>B\epsilon\)
give
\begin{equation}
0<d<\delta\alpha,\qquad 0<\alpha\le1,
\qquad d<\frac12,
\qquad 0<\omega\le\frac{\chi}{512}\le\frac1{512}.
\label{eq:q3-large-domain}
\end{equation}
In particular, \(s>\tau\), as required by Definition~\ref{def:splalm}.

\subsection{Parameters of the linear recursion}
\label{subsec:q3-large-parameters}

Set
\begin{equation}
\begin{aligned}
\nu&:=\omega^2\alpha^2,
&\zeta&:=1-\frac{\nu}{2}
       +\mathrm{i}\sqrt{\nu-\frac{\nu^2}{4}},
&u&:=\zeta-1,\\
V&:=\beta+d(1-\beta),
&E&:=\alpha^2+(1-\alpha)\nu,
&W&:=1+\frac{\chi\alpha(1-\alpha)}{E},\\
D&:=\alpha W,
&H&:=V^2+(1-V)\nu,
&\Upsilon&:=\beta V+(1-V)\nu,\\
b^2&:=\frac{\nu WH}{\eta\tau\Upsilon},
&\xi&:=\frac{\nu W(V-\beta)}{\Upsilon},
&a&:=\frac{\nu(W-1)}{\eta\alpha}+\frac{\xi}{\eta},\\
\kappa&:=\rho b^2+a,
&\vartheta&:=(1-V)\xi,
&U&:=\alpha+\chi(1-\alpha),
\end{aligned}
\label{eq:q3-large-parameters}
\end{equation}
where \(b\) denotes the positive square root. The symbol
\(\Upsilon\) is the denominator denoted by \(J\) in the proposed
construction. The radial curvature \(\kappa\) has the same meaning
as in Example~\ref{ex:rotation}; it is unrelated to the coefficient
called \(\kappa\) in~\eqref{eq:beta0-alpha}.

\begin{lemma}[Uniform parameter bounds]
\label{lem:q3-large-parameters}
\textbf{Status: proved.}
All quantities in~\eqref{eq:q3-large-parameters} are well-defined,
including when \(\beta=0\). They satisfy
\begin{equation}
\begin{gathered}
|\zeta|=1,\qquad |u|=\sqrt{\nu}=\omega\alpha,
\qquad u^2+\nu u+\nu=0,\\
\frac{\chi}{2}\le D\le1,\qquad D\ge\alpha,
\qquad 0\le\xi\le2\delta D,
\qquad 0\le\vartheta\le\delta D,\\
\frac{\omega^2\alpha^2}{\eta\tau}
\le b^2\le\frac{2\omega^2}{\eta\tau}\le\frac12,
\qquad 0\le\eta a\le2\omega^2,
\qquad 0<\kappa\le\frac12.
\end{gathered}
\label{eq:q3-large-parameter-bounds}
\end{equation}
For every \(t\in[0,1]\),
\begin{equation}
|u+t|^2=t^2+(1-t)\nu\ge\frac{\nu}{4}.
\label{eq:q3-large-denominator-bound}
\end{equation}
Finally,
\begin{equation}
\begin{aligned}
\eta a&=\frac{\nu(W-1)}{\alpha}+\xi,
&\eta\tau b^2&=\nu W+V\xi,\\
U-D&=\frac{(1-\alpha)\nu(W-1)}{\alpha},
&\beta\eta\tau b^2&=\nu(WV-\vartheta).
\end{aligned}
\label{eq:q3-large-algebra-identities}
\end{equation}
\end{lemma}

\begin{proof}
By~\eqref{eq:q3-large-domain}, \(0<\nu<1\) and \(0<V\le1\).
If \(\beta>0\), then \(\Upsilon\ge\beta V>0\).
If \(\beta=0\), then
\(V=d<1\) and \(\Upsilon=(1-d)\nu>0\).
Also \(E>0\) and \(H>0\). This proves that the definitions are valid.
The definition of \(\zeta\) gives its modulus, \(|u|^2=\nu\), and
\(u^2+\nu u+\nu=0\). Since
\(\operatorname{Re}(u)=-\nu/2\),
\[
|u+t|^2=t^2+(1-t)\nu
=(t-\nu/2)^2+\nu(1-\nu/4)
\ge\frac34\nu\ge\frac14\nu.
\]
This proves~\eqref{eq:q3-large-denominator-bound}.

Substitution of \(\nu=\omega^2\alpha^2\) gives
\[
D=\alpha+
\frac{\chi(1-\alpha)}{1+(1-\alpha)\omega^2}.
\]
It follows that
\[
D\ge\alpha,\qquad
\frac{U}{1+\omega^2}\le D\le U\le1,
\qquad U\ge\chi.
\]
Since \(\omega\le1\), these inequalities imply \(D\ge\chi/2\).

For \(\beta<1\), use
\(V-\beta=d(1-\beta)\),
\(1-V=(1-d)(1-\beta)\), and
\(\Upsilon\ge(1-V)\nu\) to obtain
\[
\frac{\xi}{W}
=\frac{\nu(V-\beta)}{\Upsilon}
\le\frac{d}{1-d}\le2d,
\]
\[
\vartheta
=\frac{\nu W(V-\beta)(1-V)}{\Upsilon}
\le W(V-\beta)\le Wd\le\delta D.
\]
Thus \(\xi\le2dW\le2\delta D\).
When \(\beta=1\), both \(\xi\) and \(\vartheta\) are zero, so the
same bounds hold.

The identity \(H-\Upsilon=V(V-\beta)\) gives
\[
\eta\tau b^2
=\frac{\nu WH}{\Upsilon}
=\nu W+V\xi.
\]
Since \(H\ge\Upsilon\) and \(W\ge1\),
\[
b^2\ge\frac{\nu}{\eta\tau}.
\]
For the upper bound, use \(\nu W=\omega^2\alpha D\),
\(V\le1\), and \(D\le1\):
\[
\eta\tau b^2
\le\omega^2\alpha D+2\delta D
\le\omega^2+2\delta\le2\omega^2.
\]
Moreover,
\[
\eta a
=\omega^2\alpha(W-1)+\xi
=\omega^2(D-\alpha)+\xi
\le\omega^2D+2\delta D\le2\omega^2.
\]
The choices of \(\omega\) in~\eqref{eq:q3-large-fixed-constants}
therefore imply
\[
b^2\le\frac{2\omega^2}{\eta\tau}\le\frac12,
\qquad
0<\kappa=\rho b^2+a
\le\frac{2\omega^2}{\eta}\left(1+\frac{\rho}{\tau}\right)
\le\frac12.
\]
Strict positivity of \(\kappa\) follows from \(\rho>0\) and \(b>0\).

The first identity in~\eqref{eq:q3-large-algebra-identities} is the
definition of \(a\); the second was proved above. For the third,
\[
U-D
=\chi(1-\alpha)\left(1-\frac{\alpha^2}{E}\right)
=\frac{\chi(1-\alpha)^2\nu}{E}
=\frac{(1-\alpha)\nu(W-1)}{\alpha}.
\]
For the last identity, the definition of \(\xi\) gives
\(\xi[\beta V+(1-V)\nu]=\nu W(V-\beta)\). Hence
\[
\begin{aligned}
\beta(\nu W+V\xi)-\nu(WV-\vartheta)
&=-\nu W(V-\beta)+\xi[\beta V+(1-V)\nu]\\
&=0.
\end{aligned}
\]
\end{proof}

\subsection{Step 1: explicit zero-initialization decomposition}
\label{subsec:q3-large-decomposition}

Define
\begin{equation}
\begin{aligned}
L&:=\frac{\tau b\zeta(u+\beta)}{u(u+V)},
&X&:=\frac{A_0}{\max\{1,|L|\}},
&\Lambda&:=LX,\\
Z&:=\frac{\alpha\zeta}{u+\alpha}X,
&M&:=\frac{\beta\zeta}{u+\beta}\Lambda,
&\bv&:=(X,\Lambda,Z,M)^{\mathsf{T}}.
\end{aligned}
\label{eq:q3-large-eigenvector}
\end{equation}
Thus \(X>0\), \(L\ne0\), and
\(\max\{X,|\Lambda|\}=A_0\).
Set
\begin{equation}
\begin{aligned}
r_x&:=-\frac{
 (W-1)\dfrac{u}{u+\alpha}
 +\xi\dfrac{\zeta}{u+V}}{W+\xi}\,X,\\
\Omega_x&:=\xi\left((1-\alpha)r_x
                  +(V-\alpha)\frac{\zeta}{u+V}X\right),\\
r_\lambda&:=\frac{\chi(X-Z)+(\eta a+D)r_x-\Omega_x}{\eta b},\\
r_z&:=r_x+X-Z,
\qquad r_\mu:=r_\lambda+\Lambda-M,\\
x_c&:=-(X+r_x),
\qquad\lambda_c:=-(\Lambda+r_\lambda),\\
\br&:=(r_x,r_\lambda,r_z,r_\mu)^{\mathsf{T}},
\qquad
\bw_c:=(x_c,\lambda_c,x_c,\lambda_c)^{\mathsf{T}}.
\end{aligned}
\label{eq:q3-large-initial-shift}
\end{equation}
All denominators are nonzero by Lemma~\ref{lem:q3-large-parameters}.
The real matrix of the shifted four-step recursion is
\begin{equation}
\mathbb{T}:=\begin{pmatrix}
1+\eta a-\chi & -\eta b & \chi & 0\\
\tau b(1+\eta a-\chi) & 1-d-\eta\tau b^2 & \tau b\chi & d\\
\alpha(1+\eta a-\chi) & -\alpha\eta b & 1-\alpha+\alpha\chi & 0\\
\beta\tau b(1+\eta a-\chi) & \beta(1-d-\eta\tau b^2)
  & \beta\tau b\chi & 1-\beta+\beta d
\end{pmatrix}.
\label{eq:q3-large-matrix}
\end{equation}
Introduce
\begin{equation}
Q(t):=(t-1+D)(t-1+V)-\vartheta(t-1+\alpha).
\label{eq:q3-large-Q}
\end{equation}

\begin{lemma}[Spectral identities and the prescribed initialization]
\label{lem:q3-large-decomposition}
\textbf{Status: proved.}
For every parameter choice above, including \(\beta=0\),
\begin{equation}
\mathbb{T}\bv=\zeta\bv,
\qquad Q(\mathbb{T})\br=0,
\qquad \bv+\br=-\bw_c.
\label{eq:q3-large-decomposition-identities}
\end{equation}
Consequently, the candidate sequence
\begin{equation}
\widetilde{\bw}^{\,k}:=
\bw_c+\zeta^k\bv+\mathbb{T}^k\br
\label{eq:q3-large-candidate}
\end{equation}
satisfies \(\widetilde{\bw}^{\,0}=0\) and
\[
\widetilde{\bw}^{\,k+1}-\bw_c
=\mathbb{T}(\widetilde{\bw}^{\,k}-\bw_c).
\]
The subspace \(\ker Q(\mathbb{T})\) is invariant under \(\mathbb{T}\).
No assertion of strict stability at \(\beta=0\) is made here.
\end{lemma}

\begin{proof}
We first prove the characteristic factorization. In
\((1+t)I-\mathbb{T}\), replace the third row by the third row minus
\(\alpha\) times the first row, and the fourth row by the fourth row
minus \(\beta\) times the original second row. Then replace the second
row by the second row minus \(\tau b\) times the first row.
These determinant-preserving operations give the matrix
\[
\begin{pmatrix}
t-\eta a+\chi&\eta b&-\chi&0\\
-\tau b(1+t)&t+d&0&-d\\
-\alpha(1+t)&0&t+\alpha&0\\
0&-\beta(1+t)&0&t+\beta
\end{pmatrix}.
\]
For a spectral variable \(1+t\), elimination of the two center
coordinates therefore gives
\[
z=\frac{\alpha(1+t)}{t+\alpha}x,
\qquad
\mu=\frac{\beta(1+t)}{t+\beta}\lambda.
\]
Substitution into the primal and dual rows produces the two equations
\[
\bigl[t(t+U)-\eta a(t+\alpha)\bigr]x
       +\eta b(t+\alpha)\lambda=0,
\]
\[
t(t+V)\lambda-\tau b(1+t)(t+\beta)x=0.
\]
The center block has determinant \((t+\alpha)(t+\beta)\),
which cancels the denominators introduced in these two equations.
Thus, initially away from these denominators and then as a polynomial
identity,
\begin{equation}
\begin{aligned}
P(t):=\det((1+t)I-\mathbb{T})
={}&t^2(t+U)(t+V)-\eta a\,t(t+\alpha)(t+V)\\
&+\eta\tau b^2(t+1)(t+\alpha)(t+\beta).
\end{aligned}
\label{eq:q3-large-characteristic-unfactored}
\end{equation}
Write
\[
A:=\frac{\nu(W-1)}{\alpha},\qquad
C:=D+V-\vartheta,\qquad N:=DV-\alpha\vartheta.
\]
Then~\eqref{eq:q3-large-algebra-identities} reads
\[
\eta a=A+\xi,\quad
\eta\tau b^2=\nu W+V\xi,\quad
U=D+(1-\alpha)A,\quad
\alpha A=\nu(W-1),\quad
\beta\eta\tau b^2=\nu(WV-\vartheta).
\]
The coefficient of \(t^3\) in \(P\) is
\[
U+V-\eta a+\eta\tau b^2
=D+V-\vartheta-\alpha A+\nu W=C+\nu.
\]
The coefficient of \(t^2\) is
\[
\begin{aligned}
UV-\eta a(\alpha+V)+\eta\tau b^2(1+\alpha+\beta)
&=N+\nu(V+1)+\nu W(\alpha+\beta-V)+\beta V\xi\\
&=N+\nu(V+1)+\nu D-\nu\vartheta\\
&=N+\nu C+\nu.
\end{aligned}
\]
Here the second equality uses
\(\nu W(\beta-V)+\beta V\xi=-\nu\vartheta\), which follows
from the last identity in~\eqref{eq:q3-large-algebra-identities}.
The coefficient of \(t\) is
\[
\begin{aligned}
-\eta a\alpha V+\eta\tau b^2(\alpha+\beta+\alpha\beta)
&=-\nu V(W-1)+\nu D
   +(1+\alpha)\nu(WV-\vartheta)\\
&=\nu(C+N).
\end{aligned}
\]
The constant coefficient is
\(\alpha\beta\eta\tau b^2=\nu N\), and the leading coefficient
is one. Hence
\begin{equation}
P(t)=(t^2+\nu t+\nu)
       \bigl((t+D)(t+V)-\vartheta(t+\alpha)\bigr).
\label{eq:q3-large-characteristic-factorization}
\end{equation}

Next, \(u^2+\nu u+\nu=0\) implies \(P(u)=0\).
Divide~\eqref{eq:q3-large-characteristic-unfactored}, evaluated at
\(u\), by \(u(u+V)(u+\alpha)\). The result is
\[
\frac{u(u+U)}{u+\alpha}-\eta a
 +\frac{\eta\tau b^2\zeta(u+\beta)}{u(u+V)}=0.
\]
Since
\(X-Z=(1-\alpha)uX/(u+\alpha)\) and
\(\Lambda/X=\tau b\zeta(u+\beta)/(u(u+V))\), this equation gives
\[
uX=\eta aX-\eta b\Lambda-\chi(X-Z).
\]
This is precisely the first row of \(\mathbb{T}\bv=\zeta\bv\).
The definition of \(\Lambda\) gives
\[
u(u+V)\Lambda=\tau b\zeta(u+\beta)X.
\]
Together with
\(\Lambda-M=(1-\beta)u\Lambda/(u+\beta)\), it implies
\[
u\Lambda=\tau b\zeta X-d(\Lambda-M),
\]
which verifies the dual row after the primal row.
The identities
\[
(u+\alpha)Z=\alpha\zeta X,
\qquad (u+\beta)M=\beta\zeta\Lambda
\]
verify the two center rows. Thus \(\mathbb{T}\bv=\zeta\bv\).

We now prove the annihilator identity without assuming stability.
For this algebraic argument only, allow \(\nu\) to vary independently
in \((0,4)\), with the other formulas unchanged and with arbitrary
\(X>0\). Initially assume
\[
0<\alpha,\beta,d<1,\qquad V\ne\alpha,
\]
and that the roots \(t_1,t_2\) of \(Q\) are distinct.
Here
\[
Q(1)=N=\frac{D\beta H}{\Upsilon}>0,
\]
\[
Q(1-\alpha)=\alpha(W-1)(V-\alpha)\ne0,
\qquad Q(1-V)=\vartheta(V-\alpha)\ne0.
\]
Therefore none of the following eigenvector denominators vanishes:
\[
\mathbf{e}(t):=
\left(
1,\frac{\tau bt(t-1+\beta)}{(t-1)(t-1+V)},
\frac{\alpha t}{t-1+\alpha},
\frac{\tau b\beta t^2}{(t-1)(t-1+V)}
\right)^{\mathsf{T}}.
\]
The elimination equations above, and the factorization of \(P\), give
\(\mathbb{T}\mathbf{e}(t_j)=t_j\mathbf{e}(t_j)\).
For coefficients \(c_1,c_2\), requiring
\(\bv+\sum_jc_j\mathbf{e}(t_j)\) to have equal first and third
coordinates, and equal second and fourth coordinates, yields
\begin{equation}
\sum_{j=1}^2c_j\frac{t_j-1}{t_j-1+\alpha}
=-\frac{u}{u+\alpha}X,
\qquad
\sum_{j=1}^2c_j\frac{t_j}{t_j-1+V}
=-\frac{\zeta}{u+V}X.
\label{eq:q3-large-synchronization-system}
\end{equation}
Indeed, the first coordinate difference in \(\mathbf{e}(t)\) is
\((1-\alpha)(t-1)/(t-1+\alpha)\); the second coordinate
difference is \((1-\beta)\tau bt/(t-1+V)\). Dividing by the
nonzero factors \(1-\alpha\) and \((1-\beta)\tau b\),
respectively, gives~\eqref{eq:q3-large-synchronization-system}.

Let \(v_j=t_j-1\). These numbers satisfy
\(v_j^2+Cv_j+N=0\). At any such root,
\[
\frac{v_j}{v_j+\alpha}
=\frac{N+\alpha v_j}{\alpha^2-C\alpha+N},
\qquad
\frac{v_j+1}{v_j+V}
=\frac{N-C+V-(1-V)v_j}{V^2-CV+N}.
\]
These identities follow by multiplying by their denominators and
replacing \(v_j^2\) by \(-Cv_j-N\).
The two constant denominators are
\[
A_\alpha:=\alpha^2-C\alpha+N
=\alpha(W-1)(V-\alpha),
\]
\[
A_V:=V^2-CV+N=\vartheta(V-\alpha).
\]
With \(S_0=c_1+c_2\) and \(S_1=c_1v_1+c_2v_2\),
the synchronization system becomes
\begin{equation}
\begin{pmatrix}
N&\alpha\\ N-C+V&-(1-V)
\end{pmatrix}
\begin{pmatrix}S_0\\ S_1\end{pmatrix}
=-X\begin{pmatrix}
A_\alpha u/(u+\alpha)\\ A_V\zeta/(u+V)
\end{pmatrix}.
\label{eq:q3-large-two-by-two-system}
\end{equation}
Its determinant is
\[
-N(1-V)-\alpha(N-C+V)
=(\alpha-V)\bigl[D(1-V)+\alpha\vartheta\bigr]
=\alpha(1-V)(\alpha-V)(W+\xi)\ne0.
\]
Cramer's rule for the first component gives
\[
S_0
=-\frac{(W-1)u/(u+\alpha)+\xi\zeta/(u+V)}{W+\xi}X
=r_x.
\]
The first equation in~\eqref{eq:q3-large-two-by-two-system}
then gives
\[
S_1=-(WV-\vartheta)r_x
       -(W-1)(V-\alpha)\frac{u}{u+\alpha}X.
\]
Using
\[
(W-1)\frac{u}{u+\alpha}X
=-(W+\xi)r_x-\xi\frac{\zeta}{u+V}X,
\]
we obtain
\[
\begin{aligned}
S_1
&=\bigl[-WV+\vartheta+(V-\alpha)(W+\xi)\bigr]r_x
  +\xi(V-\alpha)\frac{\zeta}{u+V}X\\
&=-Dr_x+\Omega_x.
\end{aligned}
\]
The coefficients themselves exist explicitly:
\[
c_1=\frac{S_1-v_2S_0}{v_1-v_2},
\qquad
c_2=\frac{v_1S_0-S_1}{v_1-v_2}.
\]
For \(\br':=c_1\mathbf{e}(t_1)+c_2\mathbf{e}(t_2)\),
its first coordinate is \(r_x\), and its synchronization equations
give \(r'_z=r_x+X-Z\) and
\(r'_\mu=r'_\lambda+\Lambda-M\).
The first row of \(\mathbb{T}-I\) gives
\[
S_1=\eta ar_x-\eta br'_\lambda+\chi(X-Z).
\]
Consequently
\[
r'_\lambda
=\frac{\chi(X-Z)+(\eta a+D)r_x-\Omega_x}{\eta b}
=r_\lambda.
\]
All four coordinates of \(\br'\) therefore equal those of \(\br\).
Since each \(t_j\) is a root of \(Q\),
\(Q(\mathbb{T})\br=0\) in this nondegenerate case.

For completeness, the continuation to excluded parameter values is
purely algebraic. On
\[
0<\alpha\le1,\quad 0\le\beta\le1,\quad 0<d<1,
\quad 0<\nu<4,
\]
all entries of \(Q(\mathbb{T})\br\) are continuous: \(b>0\),
\(W+\xi>0\), \(\Upsilon>0\), and the complex denominators
\(u\), \(u+\alpha\), \(u+\beta\), and \(u+V\) do not vanish.
The nondegenerate parameters used above are dense. To see this, first
choose interior \((\alpha,\beta,d)\) with \(V\ne\alpha\) and
\(V\ne U\), which is possible arbitrarily close to any target since
\(V=\beta+d(1-\beta)\) varies strictly with \(\beta\) when \(d<1\).
For these fixed interior parameters, the discriminant
\(C^2-4N\) is a rational function of \(\nu\) and tends to
\((U-V)^2>0\) as \(\nu\downarrow0\). Its numerator is therefore
not the zero polynomial, and its zeros are isolated. Perturbing \(\nu\)
arbitrarily little avoids them. Taking limits proves the annihilator
identity at repeated roots, \(V=\alpha\), \(\alpha=1\),
\(\beta=1\), and \(\beta=0\). One then specializes back to
\(\nu=\omega^2\alpha^2\). This continuation concerns only an
identity, not strict stability.

Finally, the definitions give, coordinate by coordinate,
\[
\bv+\br
=(X+r_x,\Lambda+r_\lambda,X+r_x,\Lambda+r_\lambda)^{\mathsf{T}}
=-\bw_c.
\]
The candidate initialization and recurrence now follow from
\(\mathbb{T}\bv=\zeta\bv\). The kernel is invariant because
\(Q(\mathbb{T})\mathbb{T}=\mathbb{T}Q(\mathbb{T})\).
\end{proof}

\subsection{Step 2: uniform bounds for the remainder}
\label{subsec:q3-large-remainder}

\begin{lemma}[A uniformly bounded second-order kernel]
\label{lem:q3-large-kernel}
\textbf{Status: proved.}
Write \(Q(t)=t^2-St+P_0\), so that
\[
S=2-D-V+\vartheta,
\qquad P_0=(1-D)(1-V)+\vartheta(1-\alpha).
\]
Both roots \(t_1,t_2\) lie in the closed unit disk. Define
\[
G_{-1}:=0,
\qquad G_j:=\sum_{i=0}^jt_1^{\,j-i}t_2^{\,i},\quad j\ge0,
\]
with the same finite sum used when the roots coincide. Then
\begin{equation}
\sup_{j\ge0}|G_j|\le\frac4D.
\label{eq:q3-large-kernel-bound}
\end{equation}
Every sequence satisfying \(y_{k+2}=Sy_{k+1}-P_0y_k\) has
\begin{equation}
y_k=G_{k-1}y_1-P_0G_{k-2}y_0,\qquad k\ge1.
\label{eq:q3-large-recurrence-solution}
\end{equation}
These statements include \(\beta=0\).
\end{lemma}

\begin{proof}
Let \(h_0=(1-\delta)D\). The parameter bounds give
\[
S\ge0,\qquad
0\le P_0\le1-D+\vartheta\le1-h_0<1,
\]
\[
Q(1)=DV-\alpha\vartheta
=\frac{D\beta H}{\Upsilon}\ge0.
\]
If the roots are real, their nonnegative sum and product imply that
both are nonnegative. If one exceeded one and the other were less
than one, \(Q(1)\) would be negative. If both were at least one and
one exceeded one, their product would exceed one. Thus both roots
belong to \([0,1]\). Order them as \(t_1\ge t_2\). Then
\[
t_2\le\sqrt{P_0}\le\sqrt{1-h_0},
\qquad
|G_j|\le\sum_{i=0}^jt_2^i
\le\frac{1}{1-\sqrt{1-h_0}}
\le\frac{2}{h_0}\le\frac4D.
\]
Here \(1-\sqrt{1-h_0}=h_0/(1+\sqrt{1-h_0})\ge h_0/2\),
and \(\delta\le1/2\).
If the roots are nonreal, they are conjugate and have common modulus
\(q=\sqrt{P_0}\le\sqrt{1-h_0}<1\). Hence
\[
|G_j|\le(j+1)q^j
\le\sum_{i=0}^jq^i
\le\frac{1}{1-\sqrt{1-h_0}}\le\frac4D.
\]
This proves the kernel bound, including repeated real roots.
The kernel satisfies \(G_0=1\), \(G_1=S\), and
\(G_{j+1}=SG_j-P_0G_{j-1}\).
Thus the right-hand side of~\eqref{eq:q3-large-recurrence-solution}
has initial values \(y_1\), \(Sy_1-P_0y_0\), and satisfies the
same recurrence. Induction proves the formula.
\end{proof}

\begin{lemma}[Remainder bounds independent of \(s\)]
\label{lem:q3-large-remainder}
\textbf{Status: proved.}
Write
\[
(p_k,\ell_k,z_k^{\rm r},m_k)^{\mathsf{T}}:=\mathbb{T}^k\br.
\]
Then, for every \(k\ge0\),
\begin{equation}
\begin{aligned}
|p_k|&\le\frac{36\omega}{\chi}X<\frac X8,
&|z_k^{\rm r}|&\le\frac{36\omega}{\chi}X,\\
|\ell_k|&\le C_\lambda|\Lambda|,
&|m_k|&\le(1+C_\lambda)|\Lambda|.
\end{aligned}
\label{eq:q3-large-remainder-bounds}
\end{equation}
In particular, \(\mathbb{T}^k\br\) is uniformly bounded also when
\(\beta=0\). None of the constants in these bounds grows with \(s\).
\end{lemma}

\begin{proof}
From the formula for \(r_x\), the inequalities
\(|u+\alpha|\ge\alpha\),
\(\xi/W\le2d\), and
\(|u+V|\ge\sqrt{\nu}/2\) give
\[
\begin{aligned}
\frac{|r_x|}{X}
&\le\frac{W-1}{W+\xi}\frac{|u|}{|u+\alpha|}
   +\frac{\xi}{W+\xi}\frac{1}{|u+V|}\\
&\le\omega+\frac{4d}{\sqrt{\nu}}
\le\omega+\frac{4\delta}{\omega}
=\frac{17}{16}\omega\le2\omega.
\end{aligned}
\]
The estimates \(\xi\le2\delta D\),
\(|V-\alpha|\le V+\alpha\),
\(V/|u+V|\le1\), and
\(\alpha/|u+V|\le2/\omega\) imply
\[
\begin{aligned}
\frac{|\Omega_x|}{DX}
&\le2\delta\left(2\omega+1+\frac{2}{\omega}\right)\\
&=4\delta\omega+2\delta+\frac{4\delta}{\omega}
=\frac{\omega^3}{16}+\frac{\omega^2}{32}+\frac{\omega}{16}
\le\frac{5\omega}{32}\le\omega.
\end{aligned}
\]
By the first row of \(\mathbb{T}\) and the definition of
\(r_\lambda\),
\[
p_0=r_x,
\qquad p_1=(1-D)r_x+\Omega_x.
\]
The identity \(Q(\mathbb{T})\br=0\) shows that \(p_k\)
satisfies the recurrence in Lemma~\ref{lem:q3-large-kernel}.
For \(k\ge1\), its solution formula therefore yields
\[
\begin{aligned}
|p_k|
&\le\frac4D\bigl(|p_1|+P_0|p_0|\bigr)\\
&\le\frac4D\bigl(2|r_x|+|\Omega_x|\bigr)\\
&\le\left(\frac{16\omega}{D}+4\omega\right)X
\le\frac{36\omega}{\chi}X.
\end{aligned}
\]
The bound also holds for \(k=0\).
Because \(\omega\le\chi/512\),
\(36\omega/\chi\le36/512<1/8\).

The center identity gives
\[
X-Z=\frac{(1-\alpha)u}{u+\alpha}X,
\qquad |X-Z|\le\omega X.
\]
Consequently \(|r_z|\le3\omega X\).
The third row of the recursion is
\[
z_{k+1}^{\rm r}=(1-\alpha)z_k^{\rm r}+\alpha p_{k+1}.
\]
Since \(0<\alpha\le1\), induction bounds \(|z_k^{\rm r}|\)
by \(\max\{|r_z|,\sup_j|p_j|\}\), which is at most
\(36\omega X/\chi\).

To estimate the multiplier without a factor proportional to \(s\),
use the primal row itself:
\[
\eta b\ell_k=(1+\eta a-\chi)p_k+\chi z_k^{\rm r}-p_{k+1}.
\]
Its first coefficient is nonnegative because \(\chi\le1\) and
\(a\ge0\). Since \(\eta a\le2\omega^2\le2\),
\begin{equation}
\eta b|\ell_k|
\le(2+\eta a)\frac{36\omega}{\chi}X
\le\frac{144\omega}{\chi}X.
\label{eq:q3-large-dual-from-primal}
\end{equation}
Let \(E_\beta:=|u+\beta|^2=\beta^2+(1-\beta)\nu\).
Using \(\operatorname{Re}(u)=-\nu/2\), we have
\[
\operatorname{Re}\bigl((u+V)\overline{(u+\beta)}\bigr)
=\Upsilon+\frac{\nu}{2}(V-\beta)
\ge\Upsilon.
\]
The modulus of this product is \(\sqrt{HE_\beta}\), so
\(\sqrt{HE_\beta}\ge\Upsilon\). The definition of \(\Lambda\)
then implies
\begin{equation}
\begin{aligned}
\eta b|\Lambda|
&=\frac{\eta\tau b^2\sqrt{E_\beta}}{\sqrt{\nu H}}X\\
&=\frac{\sqrt{\nu}W\sqrt{HE_\beta}}{\Upsilon}X
\ge\sqrt{\nu}WX=\omega DX.
\end{aligned}
\label{eq:q3-large-dual-amplitude-lower}
\end{equation}
Combining~\eqref{eq:q3-large-dual-from-primal}
and~\eqref{eq:q3-large-dual-amplitude-lower} gives
\[
|\ell_k|
\le\frac{144}{\chi D}|\Lambda|
\le\frac{288}{\chi^2}|\Lambda|
\le C_\lambda|\Lambda|.
\]

Finally,
\[
\Lambda-M=\frac{(1-\beta)u}{u+\beta}\Lambda.
\]
Since \(E_\beta\ge(1-\beta)\nu\),
\[
|\Lambda-M|^2
=\frac{(1-\beta)^2\nu}{E_\beta}|\Lambda|^2
\le(1-\beta)|\Lambda|^2\le|\Lambda|^2.
\]
Thus \(|r_\mu|\le(1+C_\lambda)|\Lambda|\).
The fourth row reads
\[
m_{k+1}=(1-\beta)m_k+\beta\ell_{k+1}.
\]
Induction gives the asserted bound for \(m_k\), including
\(\beta=0\), when \(m_k=r_\mu\) for every \(k\).
\end{proof}

\subsection{Step 3: membership in the normalized problem class}
\label{subsec:q3-large-instance}

Set \(R_*:=5X/4\), and define the instance by
\begin{equation}
\calA(x):=b(x-x_c),\qquad
h(y):=\|y\|_2,\qquad
f(x):=\psi(\|x-x_c\|)-b\langle\lambda_c,x-x_c\rangle,
\label{eq:q3-large-instance}
\end{equation}
where
\begin{equation}
\psi(r):=
\begin{cases}
-\dfrac{\kappa}{2}r^2,&0\le r\le R_*,\\[3pt]
\dfrac{\kappa}{2}(r-2R_*)^2-\kappa R_*^2,
   &R_*\le r\le3R_*,\\[3pt]
\kappa R_*(r-3R_*)-\dfrac{\kappa}{2}R_*^2,
   &r\ge3R_*.
\end{cases}
\label{eq:q3-large-radial-extension}
\end{equation}

\begin{lemma}[Uniform global constants and the initial gap]
\label{lem:q3-large-instance}
\textbf{Status: proved.}
The instance in~\eqref{eq:q3-large-instance} belongs to
\(\mathcal I(\mathbf C_0,\Delta)\). More explicitly, it satisfies
\[
L_f\le\frac12,\qquad G_f\le\frac9{32}<1,
\qquad L_{\calA}=0,\qquad G_{\calA}=b\le1,
\qquad L_h=1,
\]
and \(F(0)-\inf F\le\Delta\).
\end{lemma}

\begin{proof}
By the remainder bounds at \(k=0\),
\[
|x_c|\le X+|r_x|\le\frac98X<R_*,
\]
\[
|\lambda_c|\le|\Lambda|+|r_\lambda|
\le(1+C_\lambda)A_0\le\frac18.
\]
The last inequality follows from
\(A_0\le[8(1+C_\lambda)+5\rho]^{-1}\).
Also \(R_*\le5A_0/4\).

The values and first derivatives of \(\psi\) agree at both junctions.
At \(R_*\) the two values equal \(-\kappa R_*^2/2\), and both
first derivatives equal \(-\kappa R_*\). At \(3R_*\) both values
again equal \(-\kappa R_*^2/2\), and both first derivatives equal
\(\kappa R_*\). On the three open pieces,
\[
\psi'(r)=
\begin{cases}
-\kappa r,&r<R_*,\\
\kappa(r-2R_*),&R_*<r<3R_*,\\
\kappa R_*,&r>3R_*,
\end{cases}
\qquad
\psi''(r)\in\{-\kappa,\kappa,0\}.
\]
In particular, \(|\psi'(r)|\le\kappa R_*\).
For \(r>0\), the radial and tangential Hessian eigenvalues of
\(\psi(\|x-x_c\|)\) are \(\psi''(r)\) and \(\psi'(r)/r\).
The latter equals \(-\kappa\) in the first piece, belongs to
\([-\kappa,\kappa/3]\) in the middle piece, and belongs to
\([0,\kappa/3]\) in the last piece. Near the origin of these radial
coordinates, the function is the quadratic
\(-\kappa\|x-x_c\|^2/2\).
Its gradient is therefore continuous everywhere. Integrating its
piecewise bounded Hessian along a line segment proves that this gradient
is globally \(\kappa\)-Lipschitz. The added linear term does not change
this Lipschitz constant. Hence \(L_f\le\kappa\le1/2\).
Furthermore,
\[
\|\nabla f(x)\|
\le\kappa R_*+b|\lambda_c|
\le\frac58A_0+\frac18
\le\frac9{32}.
\]
The affine map has constant Jacobian \(bI_2\), so
\(L_{\calA}=0\) and \(G_{\calA}=b\le1\).
The Euclidean norm is convex, and the reverse triangle inequality
shows that its Lipschitz constant is one.

The radial extension satisfies \(\psi(r)\ge-\kappa R_*^2\).
Writing \(r=\|x-x_c\|\), the objective satisfies
\[
F(x)=\psi(r)-b\langle\lambda_c,x-x_c\rangle+br
\ge-\kappa R_*^2+b(1-|\lambda_c|)r
\ge-\kappa R_*^2.
\]
At zero, \(|x_c|<R_*\), so \(\psi(|x_c|)\le0\). Therefore
\[
F(0)\le b(1+|\lambda_c|)|x_c|
\le\frac98 bR_*\le\frac{45}{32}A_0.
\]
It follows that
\[
\begin{aligned}
F(0)-\inf F
&\le\frac{45}{32}A_0+\kappa R_*^2\\
&\le\frac{45}{32}A_0+\frac{25}{32}A_0^2\\
&\le\frac{205}{128}A_0
<2A_0\le\frac{\Delta}{2}\le\Delta.
\end{aligned}
\]
This verifies every constant in \(\mathbf C_0\) and the gap budget.
\end{proof}

\subsection{Step 4: both proximal outputs and the actual trajectory}
\label{subsec:q3-large-realization}

\begin{lemma}[The candidate sequence is the actual algorithm sequence]
\label{lem:q3-large-realization}
\textbf{Status: proved.}
For the instance above, the actual SPLALM iterates are
\[
\bw^k=\bw_c+\zeta^k\bv+\mathbb{T}^k\br,
\qquad k\ge0.
\]
For every \(k\ge0\), both proximal evaluations, at
\((x^k,\lambda^k)\) and \((x^{k+1},\lambda^k)\), return zero.
Moreover,
\begin{equation}
\frac34X\le|x^k-x_c|\le\frac54X=R_*.
\label{eq:q3-large-annulus}
\end{equation}
\end{lemma}

\begin{proof}
First consider only the candidate sequence
\(\widetilde{\bw}^{\,k}\) in~\eqref{eq:q3-large-candidate}.
Its first coordinate satisfies
\[
\widetilde x^{\,k}-x_c=\zeta^kX+p_k.
\]
Since \(|\zeta|=1\) and \(|p_k|<X/8\),
\[
\frac78X\le|\widetilde x^{\,k}-x_c|\le\frac98X.
\]
These bounds imply the weaker annulus in
\eqref{eq:q3-large-annulus} for every candidate point.
Its multiplier satisfies
\[
|\widetilde\lambda^{\,k}|
\le|\lambda_c|+|\Lambda|+|\ell_k|
\le2(1+C_\lambda)A_0.
\]
Consequently, for both \(j=k\) and \(j=k+1\),
\begin{equation}
\begin{aligned}
|\widetilde\lambda^{\,k}+\rho\calA(\widetilde x^{\,j})|
&\le2(1+C_\lambda)A_0+\frac54\rho bX\\
&\le\left[2(1+C_\lambda)+\frac54\rho\right]A_0
\le\frac14<1.
\end{aligned}
\label{eq:q3-large-prox-interior}
\end{equation}
The final inequality uses the definition of \(A_0\), whose third
upper bound is the reciprocal of four times the bracket.

For \(h(y)=\|y\|_2\),
\(\partial h(0)=\{q:\|q\|_2\le1\}\).
The optimality condition for the strongly convex proximal problem is
\[
y=\prox_{h/\rho}(\calA(x)+\lambda/\rho)
\quad\Longleftrightarrow\quad
\lambda+\rho(\calA(x)-y)\in\partial h(y).
\]
Thus \(|\lambda+\rho\calA(x)|\le1\) makes its unique minimizer
\(y=0\). Equation~\eqref{eq:q3-large-prox-interior} proves this
for both candidate evaluations, at every iteration.

On the candidate trajectory, all primal points lie in the quadratic
piece. With \(p=x-x_c\) and \(\ell=\lambda-\lambda_c\),
the reduced function is exactly
\[
\calD_\rho(x,\lambda)
=-\frac{\kappa-\rho b^2}{2}|p|^2+b\langle\ell,p\rangle
=-\frac a2|p|^2+b\langle\ell,p\rangle.
\]
Here \(h(0)=0\). Its gradients at the relevant evaluation points are
\[
\nabla_x\calD_\rho=-ap+b\ell,
\qquad \nabla_\lambda\calD_\rho=bp.
\]
After shifting both centers by \((x_c,\lambda_c)\), the actual four
updates consequently read
\[
\begin{aligned}
p^+&=(1+\eta a-\chi)p-\eta b\ell+\chi z,\\
\ell^+&=\ell+\tau bp^+-d(\ell-m),\\
z^+&=(1-\alpha)z+\alpha p^+,\\
m^+&=(1-\beta)m+\beta\ell^+.
\end{aligned}
\]
These are precisely the rows of~\eqref{eq:q3-large-matrix}.
The candidate sequence satisfies this recurrence and starts from zero
by Lemma~\ref{lem:q3-large-decomposition}. It also satisfies every
proximal and quadratic-region condition just verified.
Therefore, at iteration zero its four updates equal the actual updates.
If the two sequences agree through iteration \(k\), the unique primal
update and the unique proximal output at the new primal point agree,
and then the dual and center updates agree as well. Induction proves
identity of the two sequences for all \(k\).
\end{proof}

\subsection{Step 5: a residual lower bound at every iteration}
\label{subsec:q3-large-residual}

\begin{lemma}[Uniform residual lower bound]
\label{lem:q3-large-residual}
\textbf{Status: proved.}
For every \(k\ge0\), the actual post-primal output satisfies
\[
R_k\ge\frac34 bX\ge c_{\rm r}\alpha.
\]
\end{lemma}

\begin{proof}
The previous lemma gives \(y^{k+1}=0\), and hence
\[
q^k=\lambda^k+\rho\calA(x^{k+1})\in\partial h(0).
\]
The prescribed KKT residual is therefore
\[
R_k=\max\left\{
\left|-a(x^{k+1}-x_c)+b(\lambda^k-\lambda_c)\right|,
\ b|x^{k+1}-x_c|\right\}.
\]
Using only its feasibility component and the annulus gives
\(R_k\ge3bX/4\).

It remains to lower-bound the normalized amplitude \(bX\).
The parameter bounds give
\[
b\ge\frac{\sqrt{\nu}}{\sqrt{\eta\tau}}
=\frac{\omega\alpha}{\sqrt{\eta\tau}}.
\]
Also, since \(|u+V|=\sqrt H\) and
\(|u+\beta|=\sqrt{E_\beta}\),
\[
\sqrt{E_\beta}=|u+\beta|
\le|u+V|+(V-\beta)\le\sqrt H+d.
\]
By \(\sqrt H\ge\sqrt\nu/2\) and \(d\le\delta\alpha\),
\[
\frac{d}{\sqrt H}
\le\frac{2\delta}{\omega}=\frac{\omega}{32}\le1.
\]
Consequently \(\sqrt{E_\beta}\le2\sqrt H\).
The definition of \(L\) now yields
\[
\frac{b}{|L|}
=\frac{\sqrt\nu\sqrt H}{\tau\sqrt{E_\beta}}
\ge\frac{\omega\alpha}{2\tau}.
\]
Since \(L\ne0\) and \(X=A_0/\max\{1,|L|\}\),
\[
\begin{aligned}
bX
&=A_0\min\left\{b,\frac{b}{|L|}\right\}\\
&\ge A_0\omega\alpha
\min\left\{\frac{1}{\sqrt{\eta\tau}},\frac{1}{2\tau}\right\}.
\end{aligned}
\]
Multiplication by \(3/4\) gives the asserted constant
\(c_{\rm r}\) in~\eqref{eq:q3-large-fixed-constants}.
\end{proof}

\begin{proof}[Proof of Theorem~\ref{thm:q3-large-alpha}]
The conditions \(0<\epsilon<\epsilon_0\) and
\(\alpha>B\epsilon\) imply~\eqref{eq:q3-large-domain}.
Lemmas~\ref{lem:q3-large-parameters}--\ref{lem:q3-large-decomposition}
therefore define the instance parameters and its zero-initialized
candidate trajectory. Lemma~\ref{lem:q3-large-instance} places the
instance in \(\mathcal I(\mathbf C_0,\Delta)\), and
Lemma~\ref{lem:q3-large-realization} proves that this is the actual
algorithm trajectory. By Lemma~\ref{lem:q3-large-residual} and
\(B\ge2/c_{\rm r}\),
\[
R_k\ge c_{\rm r}\alpha
>c_{\rm r}B\epsilon\ge2\epsilon>\epsilon
\qquad\text{for every }k\ge0.
\]
No output satisfies the tolerance. Definition~\ref{def:complexity-measure}
therefore gives \(T^{\mathrm{KKT}}_\epsilon=+\infty\).
\end{proof}

\subsection{Step 6: separation from the proved upper-bound path}
\label{subsec:q3-large-upper-separation}

\begin{proposition}[Nonoverlap with Theorem~\ref{thm:upper-beta0}]
\label{prop:q3-large-upper-separation}
\textbf{Status: proved.}
Whenever the fixed parameters also satisfy the hypotheses of
Theorem~\ref{thm:upper-beta0}, its prescribed stepsize
\[
\alpha_{\rm ub}=\frac{\kappa_{\rm ub}}{\ell s},\qquad\beta=0,
\]
where \(\kappa_{\rm ub}\) denotes the coefficient in
\eqref{eq:beta0-alpha}, satisfies
\[
\alpha_{\rm ub}<\frac{d}{384}<B\epsilon
\]
for all sufficiently small \(\epsilon\) at the same
\(s=c_s/\epsilon\). Thus the two conclusions have disjoint
parameter ranges.
\end{proposition}

\begin{proof}
Use the upper-bound appendix's notation, with subscripts here to avoid
confusing it with the present instance parameters:
\[
r_{\rm ub}:=\gamma^{-1},\qquad
b_s^{\rm ub}:=\ell+s^{-1},\qquad
 a_s^{\rm ub}:=r_{\rm ub}-b_s^{\rm ub}>0.
\]
The quantity in~\eqref{eq:beta0-stepsizes} is
\[
(\omega_s^{\rm ub})^2
=s\frac{[a_s^{\rm ub}
 +\tau b_s^{\rm ub}(3r_{\rm ub}-2b_s^{\rm ub})]^2}
 {\tau^2(a_s^{\rm ub})^3}.
\]
Since \(b_s^{\rm ub}<r_{\rm ub}\),
\(3r_{\rm ub}-2b_s^{\rm ub}>0\); hence
\[
(\omega_s^{\rm ub})^2\ge\frac{s}{\tau^2a_s^{\rm ub}}.
\]
The third upper bound on \(\alpha\) in
\eqref{eq:beta0-stepsizes} consequently gives
\[
\alpha_{\rm ub}
\le\frac{1}{384r_{\rm ub}\tau(\omega_s^{\rm ub})^2}
\le\frac{\tau a_s^{\rm ub}}{384r_{\rm ub}s}
<\frac{\tau}{384s}=\frac{d}{384}.
\]
On the other hand, our threshold satisfies
\[
B\epsilon\ge\frac{\tau\epsilon}{c_s\delta}
=\frac{d}{\delta}>\frac{d}{384},
\]
because \(0<\delta<1\).
This proves, in particular, the requested strict inequality
\[
B\epsilon>\frac{\kappa_{\rm ub}}{\ell s}.
\]
No extra dependence of \(B\) on an unspecified upper-bound coefficient
is needed.
\end{proof}

\subsection{The difference between \(\beta=0\) and \(\beta>0\)}
\label{subsec:q3-large-beta-endpoint}

\paragraph{Status: proved.}
The construction, all six steps, and their constants are uniform over
\(\beta\in[0,1]\). Their spectral interpretations differ as follows.

For \(\beta>0\),
\[
Q(1)=\frac{D\beta H}{\Upsilon}>0.
\]
The root argument in Lemma~\ref{lem:q3-large-kernel} then puts both
roots strictly inside the unit disk: real roots lie in \([0,1]\),
and the displayed strict inequality excludes one; nonreal roots have
modulus \(\sqrt{P_0}<1\).
Each component of \(\mathbb{T}^k\br\) satisfies the scalar
second-order recurrence. For distinct roots, its solution is a linear
combination of their powers. For a repeated root \(t\), the solution
has the form \((c_0+c_1k)t^k\), with the zero-root case interpreted
directly from the recurrence. These expressions tend to zero.
Thus the remainder decays for each fixed \(\beta>0\), although no
uniform decay rate as \(\beta\downarrow0\) is needed or asserted.

For \(\beta=0\), the parameter formulas instead give exactly
\[
V=d,\qquad \Upsilon=(1-d)\nu,
\qquad \xi=\frac{Wd}{1-d},\qquad \vartheta=Wd.
\]
Therefore
\begin{equation}
Q(t)=(t-1)(t-r_0),
\qquad r_0:=1-D+d(W-1).
\label{eq:q3-large-beta-zero-factorization}
\end{equation}
Because \(W\ge1\), \(D\le1\), and \(Wd\le\delta D\),
\[
0\le1-D\le r_0
\le1-(1-\delta)D<1,
\qquad
1-r_0\ge(1-\delta)D\ge\frac{\chi}{4}.
\]
The unit root is simple. There is no Jordan growth at that root on
\(\ker Q(\mathbb{T})\), because the two factors in
\eqref{eq:q3-large-beta-zero-factorization} are distinct.
More explicitly, define
\[
\br^{(1)}:=\frac{(\mathbb{T}-r_0I)\br}{1-r_0},
\qquad
\br^{(0)}:=\frac{(I-\mathbb{T})\br}{1-r_0}.
\]
The annihilator identity proves directly that
\[
\br=\br^{(1)}+\br^{(0)},\qquad
\mathbb{T}\br^{(1)}=\br^{(1)},\qquad
\mathbb{T}\br^{(0)}=r_0\br^{(0)}.
\]
Consequently,
\begin{equation}
\mathbb{T}^k\br=\br^{(1)}+r_0^k\br^{(0)},\qquad k\ge0.
\label{eq:q3-large-beta-zero-remainder}
\end{equation}
The denominator \(1-r_0\) is bounded below by \(\chi/4\).
Independently of this decomposition, Lemma~\ref{lem:q3-large-remainder}
already provides the stated componentwise uniform bounds, so no
summation of a nondecaying remainder is used.

In this endpoint case \(M=0\), and the fourth row of \(\mathbb{T}\)
is the identity on the fourth coordinate. Moreover,
\[
r_\mu=r_\lambda+\Lambda=-\lambda_c.
\]
Thus the actual dual center is exactly
\[
\mu^k=\lambda_c+\zeta^kM+(\mathbb{T}^k\br)_\mu
=\lambda_c+r_\mu=0\qquad(k\ge0),
\]
as required by the algorithm with \(\beta=0\) and zero initialization.
The persistent component in~\eqref{eq:q3-large-beta-zero-remainder}
need not vanish; its primal component is nevertheless bounded together
with the entire remainder by \(X/8\). It therefore cannot cancel the
unit-modulus rotational component of primal amplitude \(X\).
This is the precise replacement for the strict-stability argument used
in the original \(\beta>0\) rotation branch.
